\documentclass[anon,12pt]{clear2025} % Anonymized submission
% \documentclass[final,12pt]{clear2025} % Final submission
% \documentclass[final,12pt,noproceedings]{clear2025} % Final submission for those who opt out of the proceedings; they should use clear2025-extended-abstract.tex for their extended abstract

% --- PACKAGES ---
\newtheorem{observation}[theorem]{Observation}
\usepackage[capitalise]{cleveref}
\usepackage{amsfonts}
\usepackage{bbm}

\usepackage{tikz}
\usetikzlibrary{bayesnet}

\usetikzlibrary{shapes,decorations,arrows,calc,arrows.meta,fit,positioning}
\tikzset{
    -Latex,auto,node distance =0.5 cm and 0.5 cm,semithick,
    state/.style ={ellipse, draw, minimum width = 0.7 cm},
    point/.style = {circle, draw, inner sep=0.04cm,fill,node contents={}},
    bidirected/.style={Latex-Latex,dashed},
    el/.style = {inner sep=2pt, align=left, sloped}
}


\usepackage{url}
\usepackage{booktabs}

\tikzset{%
  zeroarrow/.style = {-stealth,dashed},
  onearrow/.style = {-stealth,solid},
  c/.style = {circle,draw,solid,minimum width=2em,
        minimum height=2em},
  r/.style = {rectangle,draw,solid,minimum width=2em,
        minimum height=2em}
}

\newcommand{\midlinewidth}{1.0pt}
\newcommand{\incmidlinewidth}{1.7pt}
\newcommand{\middist}{20pt}
\newcommand{\halfdist}{13pt}
\newcommand{\pn}{\ensuremath{\mathit{PN}}}
\newcommand{\ps}{\ensuremath{\mathit{PN}}}

\definecolor{petroil2} {RGB} {36, 165, 175}
\definecolor{gold2} {RGB} {255, 130, 0}

% # CODE SNIPPET STYLE (settings borrowed from Edward ICLR paper)
\usepackage{listings}
\usepackage[dvipsnames]{xcolor}

\definecolor{strings}{rgb}{.624,.251,.259}
\definecolor{keywords}{rgb}{.224,.451,.686}
\definecolor{comment}{rgb}{.322,.451,.322}

\lstdefinelanguage{python}{
  morekeywords={from, import, as, for, in, while, def, return, 
  =, +, -, /, *, lambda},
  % keywords=[2]{build_toy_dataset, neural_network, __init__},
  keywords=[3]{sample, param, module, Marginal, Posterior, Trace, Poutine, Distribution},
  morecomment=[l]{\#},
  morecomment=[s]{"""}{"""},
  morestring=[b]',
  morestring=[b]",
  alsoletter={<>=-+/*},
  sensitive=true
}

\usepackage[most]{tcolorbox}
\newtcolorbox{docode}[1][]{
    colback=gray!20,        % background color
    colframe=black,         % border color
    width=0.9\textwidth,    % box width
    arc=4mm,                % rounded corners
    boxrule=0.5pt,          % border thickness
    left=2em,               % left padding
    enhanced,
    fontupper=\ttfamily,    % font inside the box
    #1                      % allow optional overrides
}



\lstset{
  language=python,
  keywordstyle=\color{BrickRed}\bfseries\ttfamily,
  keywordstyle=[2]\color{Violet}\ttfamily,
  keywordstyle=[3]\color{keywords}\ttfamily,
  commentstyle=\color{comment}\ttfamily,
  stringstyle=\color{strings}\ttfamily,
  basicstyle=\fontsize{8pt}{8.25pt}\selectfont\ttfamily,
  basewidth=0.5em,
  % numbers=left,
  % numberstyle=\tiny,
  % stepnumber=1,
  columns=fixed,
  xleftmargin=2ex,%0.25in,
  % firstnumber=1,
  showstringspaces=false,
  mathescape=true,
  keepspaces=True,
  tabsize=2
}
\renewcommand{\texttt}[1]{\lstinline[basicstyle=\fontsize{8pt}{8.25pt}\selectfont\ttfamily]{#1}}


\usepackage{xcolor}   % for colors
\usepackage{hyperref}

% Configure hyperlinks
\hypersetup{
    colorlinks=true,      % use colored text for links
    linkcolor=blue,       % color of internal links (sections, TOC)
    urlcolor=blue,        % color of URLs
    citecolor=blue,       % color of citations
    filecolor=blue        % color of file links
}

\title[Explainable]{A Computationally Feasible Framework for  Causal Probabilistic Explanation}
\usepackage{times}
% Use \Name{Author Name} to specify the name.
% If the surname contains spaces, enclose the surname
% in braces, e.g. \Name{John {Smith Jones}} similarly
% if the name has a "von" part, e.g \Name{Jane {de Winter}}.
% If the first letter in the forenames is a diacritic
% enclose the diacritic in braces, e.g. \Name{{\'E}louise Smith}

% Two authors with the same address
% \clearauthor{\Name{Author Name1} \Email{abc@sample.com}\and
%  \Name{Author Name2} \Email{xyz@sample.com}\\
%  \addr Address}

% Three or more authors with the same address:
% \clearauthor{\Name{Author Name1} \Email{an1@sample.com}\\
%  \Name{Author Name2} \Email{an2@sample.com}\\
%  \Name{Author Name3} \Email{an3@sample.com}\\
%  \addr Address}

% Authors with different addresses:
\clearauthor{%
 \Name{Rafal Urbaniak} \Email{rafal@basis.ai} \\
 \Name{Eli Bingham} \Email{eli@basis.ai}\\
 \addr{Basis Research Institute}\\
 \Name{Poorva Garg} \Email{poorvagarg@cs.ucla.edu}\\
 \addr{University of California, Los Angeles} \\
 \Name{Jack Feser} \Email{jack@basis.ai} \\
 \Name{Daniel Waxman} \Email{dan@basis.ai} \\
 \addr{Basis Research Institute}
}

\begin{document}

\maketitle

\begin{abstract}%
  This is a great paper and it has a concise abstract.%
\end{abstract}

\begin{keywords}%
  List of keywords%
\end{keywords}


\section{Introduction}

\texttt{ChiRho} is a causal probabilistic programming language that enables users to write causal probabilistic models as programs and compute causal queries. Interventions are performed as program transformations agnostic to the inference algorithms, and higher-level causal queries are expressed through interventions. 

Our focus in this work is answering queries related to causal explanations, i.e., given a causal model and a factual outcome, what factors were actually responsible for that outcome. The development is guided by two ideas: Pearl's probabilities of causation \cite{pearl2022probabilities}, and the Halpern-Pearl definition of causality \cite{actualCausalityHalpern}, often called \emph{actual causality}. Actual causality seeks to understand sets of causes for specific events. This is in contrast to the \emph{type causality}, more common to machine learning and statistics, which makes population-level statements of cause and effect. For example, it is now well-known that smoking cigarettes predisposes someone to lung cancer at much higher risks than the general population, i.e., smoking \emph{causes} cancer in the type-causal sense. In contrast, an actual causality asks questions such as ``did this specific lung cancer patient develop cancer as a result of their smoking?'' Answers to type-level and individual-level questions, clearly, can be different.

Unfortunately, actual causality can be somewhat restrictive in probabilistic settings, especially over continuous random variables. We therefore propose a related framework based on what we call the \emph{explanation-extended distribution}. We show how the explanation-extended distribution can be used to compute quantities of interest, including a \emph{probability of causation} and its counterpart for continuous outcomes, \emph{causal impact}. 

This paper is organized as follows. In \Cref{sec:intuitions} we motivate the developments and place them in the context of the existing related probabilistic notions. In \Cref{sec:definitions} we define the probability of causation and related notions.  \Cref{sec:relation_with_ac} relates our approach to actual causality. In \cref{sec:scaling}, we discuss scaling up the methods to realistic models and handling continuous variables. \Cref{sec:case_studies} presents case studies, including the causally aware probabilistic automated valuation model. Proofs are provided in the appendix (\cref{app:proofs}).  

%\subsection*{Notation}
%@DAN: I already introduced notation a bit later, take a look at the whole paper before you decide on writing anything here


% The rest of this work is structured as follows: in \cref{sec:intuitions}, we discuss intuitions behind our methods, which are formally defined in \cref{sec:definitions}. In \cref{sec:relation_with_ac}, we discuss our framework's relationship to actual causality. We provide a case study in \cref{sec:case_study}, and a discussion in \cref{sec:discussion}. We introduce computational methods for computing these quantities, implemented in the \texttt{Explainable} module of \texttt{ChiRho}.

\section{Motivations and intuitions}
\label{sec:intuitions}

To motivate the technical decisions and definitions that will follow, let us first think through a simplified example involving responsibility attribution. We will subsequently try to apply some of the existing ideas in causal responsibility attributions to the running example, pointing out how the example suggests  paths of improvement. We postpone the official definitions and a more technical treatment to a later section---for now, we are putting the developments in a context, and building motivations for some of the design decisions.

\begin{example}[Old Boys' Club Bank] Bob applied for a home mortgage loan with the local branch of Old Boys’ Club Bank (OBCB), and his application was denied. The OBCB policy requires applicants to pass a credit score check, but it only performs this check for male applicants; all non-male applicants are rejected outright. \label{ex:obcb}
\end{example}
\noindent Assuming this policy holds, why was Bob’s loan application denied? An intuitive explanation is that his credit score was insufficient—had his score been higher, the outcome would have been different. In contrast, his gender does not appear to be responsible: had it been different, his application still would have been denied.


% \tikz{ %
%         \node[latent] (g) {Gender} ; %
%         \node[latent, right=of g] (cs) {Credit Score} ;
%         \node[latent, left=of cs, right=of g, below=of g] (ccf) {Credit Check Failed} ;
%         \node[latent] (aa, yshift=0.5cm) {Application Accepted}
%         % \node[latent, left=of y, yshift=0.5cm] (mu) {$\mu$} ; %
%         % \node[latent, left=of y, yshift=-0.5cm] (tau) {$\tau$} ; %
%         \edge {} {} ; %
%       }

\vspace{2mm}
\noindent \textbf{But-For Causality with Single Suspects.} This suggests a counterfactual approach to feature responsibility attribution: we can identify the features responsible for a particular outcome by asking counterfactual questions—would the outcome have been different if a specific feature had a different value? This is called the \textit{but-for clause}: Bob wouldn’t have been denied the loan but for his insufficient credit score.  

To be a bit more concrete, consider a model $M$ consisting of a set of \emph{exogenous variables} $\mathbf{U}$, which are the source of stochasticity,  and \emph{endogenous variables} $\mathbf{V}$, which are influenced by other variables in the model. Let $\mathbf{S} \subseteq \mathbf{V}$ be a subset of variables we suspect to be causally responsible (call them \emph{causal suspects}), and $Y \in  \mathbf{V}$ be an outcome of interest. Suppose we  consider only one suspect $S \in \mathbf{S}$ and mark factual values with $^\star$. Say in the \emph{factual} scenario  $S = s^\star$ and $Y = y^\star$ under model $M$ and exogenous variables $\mathbf{U} = \mathbf{u}$, denoted $(M, \mathbf{u}) \models S=s^\star, Y=y^\star$. A first attempt at defining the relevant counterfactual notion of causal responsibility might be:
\begin{align*}
S=s^\star \text{ is causally responsible for } Y=y^\star \text{ in } (M, \mathbf{u}) \text{ if and only if } \\   (M, \mathbf{u})[S \leftarrow s'] \not \models Y=y^\star \text{ and } (M, \mathbf{u}) \models S=s^\star, Y=y^\star
\end{align*}
for some $s' \neq s^\star$.
The arrow represents an intervention on $S$, setting it to an alternative value. The idea is: in a given context, $S$ is causally responsible for an outcome only when intervening on $S$ with some alternative value would change that outcome. Bob’s credit score satisfies this condition: had his credit score been higher, the loan might have been approved.  

However, this notion turns out to be insufficiently fine-grained. Consider a slightly more complicated case: Alice, a woman with a credit score that would have been insufficient if checked, applies with OBCB for the same loan. Her application was rejected, but for a different reason---her gender alone was sufficient for rejection, and no credit score check was performed. Why was Alice denied the loan?  

Under the but-for clause, neither her gender nor her credit score individually appears responsible. No matter how high her credit score, her application would still have been denied because the bank does not evaluate credit scores for non-male applicants. Similarly, her gender alone does not fully explain the outcome—male applicants are denied loans based on their credit score if it is too low, so had her gender been different, her application would still have been rejected (as her credit score was insufficient anyway).

\vspace{2mm}
\noindent \textbf{But-For Causality with  Suspect Sets.} One way to address such cases is to allow sets or vectors of variables to be jointly responsible. Instead of a single suspect, consider a set of variables $\mathbf{C}\subseteq \mathbf{S}$:
\begin{align*}
\mathbf{C} = \mathbf{c}^\star \text{ is causally responsible for } Y = y^\star \text{ in } (M, \mathbf{u}) \text{ just in case } \\  (M, \mathbf{u})[\mathbf{C} \leftarrow \mathbf{c'}] \not \models Y= y ^\star
\end{align*}
for some $\mathbf{c'}$ everywhere different from $\mathbf{c}$, where no proper subset of $\mathbf{C}$ has this property--the idea here is that we want to prefer simpler explanations and not to include irrelevant properties in them. This \emph{minimality condition} looks at least like a decent first stab at an explication of this idea.


Under this generalization, the set (\texttt{gender = female, credit score = any value}) is jointly responsible for Alice's application being denied. Changing these variables together---e.g., setting (\texttt{gender = male, credit score = sufficient})---would change the decision. Thus, her denial was caused by the interplay of both features.

But what can we say about the individual features? Neither Alice’s gender nor her credit score is individually deemed responsible according to this approach. One possible reaction is to say that the whole set bears responsibility, even though the features taken separately do not. Another is to argue that any feature participating in a responsible set is partially accountable by the same token.  

Both approaches have issues, as they imply symmetry: both features are deemed equally responsible. This conflicts with the intuition that since OBCB did not even evaluate Alice’s credit score, her denial was actually caused by her gender, not by her credit score.\footnote{Examples of this sort, involving both over-determination and undercutting of causal impact, abound. Suppose Jim is both stabbed and poisoned and dies as a result. Because stabbing works quickly, it kills Jim before the poison takes effect. In this case, it’s the stabbing that caused his death. Yet analogous but-for clauses fail to break the symmetry here as well.}

\vspace{2mm}


\begin{figure} 
    \centering
    \hspace{-1cm}
    \begin{subfigure}[t]{0.42\textwidth}
    \centering
        \begin{tikzpicture}
            \node[state] (1) {\small Rejected};
            \node[state] (2) [above =of 1,xshift=1.7cm] {\small Credit Score};
            \node[state] (3) [above =of 1,xshift=-1cm] {\small Gender};
        
            \path (2) edge (1);
            \path (3) edge (1);
        \end{tikzpicture}
        \caption{Na\"ive Causal Model}
    \end{subfigure}
    \begin{subfigure}[t]{0.46\textwidth}
    \centering
        \begin{tikzpicture}
            \node[state] (1) {\small Rejected};
            \node[state] (2) [above =of 1,xshift=1.7cm] {\small Credit Check Failed};
            \node[state] (3) [above =of 1,xshift=-2cm] {\small Gender};
            \node[state] (4) [above =of 2] {\small Credit Score};
        
            \path (2) edge (1);
            \path (3) edge (1);
            \path (4) edge (2);
            \path (3) edge (2);
        \end{tikzpicture}
        \caption{Modified Causal Model}
    \end{subfigure}
    \caption{Two causal models for rejection at the OBCB. In the model in which we pay no attention to the mediator, the credit score is directly linked to potential rejection. In the modified causal model, a potential witness node is added, showing that the effect of credit score is mediated through the failure of a credit check.}
    \label{fig:obcb_causal_models}
\end{figure}

If we keep this assumption at its factual value (i.e., \texttt{CreditCheckFailed} = 0), then counterfactually, Alice would have been granted the loan if she were male. Asymmetrically, there is no context fixed at the actual value such that, had her credit score been higher, she would have been granted the loan---her gender blocked the credit check in the first place. This motivates the following formalization. Now we not only consider some causal set of suspects $\mathbf{C} \subseteq \mathbf{S}$, but also some set of context nodes $\mathbf{T} \subseteq \mathbf{W}$, where $\mathbf{W}$ is the set of all potential context nodes we could consider, also called the \emph{witness set}.
\begin{align*}
\mathbf{S} = \mathbf{s}^\star \;\text{ is causally responsible for }\; Y = y^\star \;\text{ in }\; (M, \mathbf{u}) \;\text{ just in case } \\
(M, \mathbf{u})[\mathbf{C} \leftarrow \mathbf{c}', \mathbf{T} \leftarrow \mathbf{t}^\star] \not\models Y = y,
\end{align*}
\noindent for some \(\mathbf{c}'\) distinct everywhere from  \(\mathbf{c}^\star\), and for some set of context variables \(\mathbf{T}\) (called \emph{active witnesses}) with their actual values \(\mathbf{t}^\star\), where no proper subset of \(\mathbf{C}\) satisfies this property. This notion is identical to but-for causality with suspect sets, aside from existentially quantifying over factual context $T \leftarrow \mathbf{t}^\star$ and keeping them fixed by interventions.


\vspace{2mm}
\noindent \textbf{Suspects, Witnesses, and their Activations.} This formulation explicates actual causality by evaluating but-for clauses with certain contextual assumptions held fixed. The actual causality claim holds if a context exists where the corresponding restricted but-for clause holds.   To answer causal responsibility queries using this or similar notions, we need to perform two searches: one finds causally responsible sets of variables and another determines context sets to hold fixed. The set of all candidate cause nodes will be called \emph{suspects} and denoted by $\mathbf{S}$. Subsets of $\mathbf{S}$, when considered intervened in a particular search, will be called \emph{active suspects} and denoted by $\mathbf{C}$. The candidate context variables are called \emph{witnesses} and denoted by $\mathbf{W}$. In a particular search step, some subset of $\mathbf{W}$ will be in fact intervened to be kept at the factual values; we call those \emph{active witnesses} and denote them by $\mathbf{T}$.



\vspace{2mm}
In the process of going through existing approaches, and pointing out which elements thereof we would like to preserve and why, let's move to another strategy, that of prescribing probabilities to the necessity, sufficiency, or joint necessity and sufficiency of an event.


\vspace{2mm}
\noindent \textbf{Probability of Necessity.} While we do agree that bringing witness nodes into the picture to break symmetries is a great idea, we will depart from this notion in some respects. For one thing, this definition is formulated in terms of an existential quantifier: it is enough that one actual context being kept fixed validates the counterfactual dependency. We will prefer a more subtle approach where information from counterfactual dependencies across different witness sets is used and integrated using probability distributions. After all, the mere existence of a context in which a counterfactual dependency arises might come about in importantly different ways. The context might have been a very rare combination of background factors leading to a very unlikely accidental dependency, or it might have been the presence of obvious factors everyone should expect to be present. This should make a difference to the responsibility attribution.  Moreover, we would like to be probabilistic about responsibility attribution; this can be of great importance in applications, where strict, deterministic causal models are typically not available, and our best models may be stochastic. Let us illustrate this with the following modification of \cref{ex:obcb}: 
\begin{example}[Old Boys' Club Bank, Stochastic]
Consider the Old Boys' Club Bank once again but with a stochastic policy. In this case, they proceed with credit score checks for a randomly selected 20\% of women, and do not perform credit checks for a randomly selected 10\% of men. Unchecked applicants are rejected. Furthermore, if a man's credit check fails, he is nevertheless accepted 5\% of the time; if a woman's credit check succeeds, she is denied 10\% of the time. \label{ex:obcb_stochastic}
\end{example} 

First, we consider the counterfactual impact of changing individual features. Ben is a male with insufficient credit. His probability of getting a loan is $0.9 \times 0.05 = 0.045$. If he were female, this chance would decrease to $0.2 \times 0=0$.
If he had a sufficient credit score, his probability of getting the loan would jump to $0.9$, as his only hurdle would be to obtain a credit check. Therefore, changing his gender does not increase the probability of the outcome being different (quite the opposite), while changing his credit score does significantly. Probabilistic responsibility attribution should be sensitive to such differences.

The probability discussed above is the \emph{probability of necessity} \cite{pearl2022probabilities}.
This is the probability that the outcome $Y=y$ would not have occurred in the absence of  $C = c^\star$, given that both occurred: 
\begin{align*}
\pn(C = c^\star, Y = y^\star) 
  = P\!\left(Y_{c'} \neq y^\star \;\middle|\; C = c^\star,\, Y = y^\star\right)
\end{align*}
where $Y_{c'}$ is the value of $Y$ after setting $C$ to an alternative value $c'\neq c^\star$.

While the probability of necessity is a useful step in probabilistic attribution, the notion is still insufficient for our purposes. One missing aspect is a lack of a conditional notion:
%These notions---as they are not conditional,,in a way conditional average treatment effect is a conditional version of the average treatment effect--- aren't general enough to address the queries we are interested in, as 
It is somewhat unclear how to compute $PN$ for gender and $PN$ for credit sufficiency for Ben, assuming his application was denied. But suppose we slightly generalize and compute what naturally seem to be conditional probabilities of necessity. That is, we condition both on Ben's gender and his credit insufficiency, and then consider two different interventions: on gender and on credit sufficiency, separately computing the relevant probabilities. Call counterfactual worlds in which we intervene with such alternative values the \emph{necessity worlds}. We, effectively, want to evaluate probabilities in specific necessity worlds as follows:
\begin{align*}
\mathit{PN}^{\text{ben}}(\mathtt{acc} = 0, \mathtt{gen} = 0) 
  &= P\!\left(
      \mathtt{acc}_{\mathtt{gen} = 1} \;\middle|\;
      \mathtt{gen} = 0,\, \mathtt{cred} = 0
    \right) = 0 \\
\mathit{PN}^{\text{ben}}(\mathtt{acc} = 0, \mathtt{cred} = 0) 
  &= P\!\left(
      \mathtt{acc}_{\mathtt{cred} = 1} \;\middle|\;
      \mathtt{gen} = 0,\, \mathtt{cred} = 0
    \right) = 0.9
\end{align*}
\noindent (where  $\mathtt{gen} = 0$ for males, and $\mathtt{cred} = 1$ means credit sufficiency) and $\mathtt{acc}=1$ means loan application has been accepted.


What about Alice, a female with insufficient credit? Currently, her probability of getting the loan is $0$. If she were male, her probability would rise to $0.045$, the same as Ben's. If she were to have a sufficient credit score while remaining female, the probability of a different outcome would increase to $0.2 * 0.9 = 0.18$.
\begin{align*}
\mathit{PN}^{\text{alice}}(\mathtt{acc} = 0, \mathtt{gen} = 1) 
  &= P\!\left(
      \mathtt{acc}_{\mathtt{gen} = 0} \;\middle|\;
      \mathtt{gen} = 1,\, \mathtt{cred} = 0
    \right) = 0.045 \\
\mathit{PN}^{\text{alice}}(\mathtt{acc} = 0, \mathtt{cred} = 0) 
  &= P\!\left(
      \mathtt{acc}_{\mathtt{cred} = 1} \;\middle|\;
      \mathtt{gen} = 1,\, \mathtt{cred} = 0
    \right) = 0.18
\end{align*}
Based solely on the probability of necessity, one might conclude that Alice's gender is barely responsible for the decision and that her credit score is several times less influential than Ben's. This conclusion feels misguided: while the probability of necessity properly communicates that there are many other likely scenarios that lead to Alice's application being rejected, even if she were male, it does not properly account for the fact that gender alone made the outcome extremely likely. This leads to an analogous construction of the \emph{probability of sufficiency}.




\noindent \textbf{Probability of Sufficiency.} Perhaps, to properly judge cases such as Alice's insufficient credit, we should ask how \emph{sufficient} a feature is for an outcome: how often an intervention that mimics the actual state of the candidate would cause the outcome?\footnote{Example:  the probability of necessity of breathing for doing a PhD in computer science is very high, but the corresponding probability of sufficiency is rather low.}   We will call possible worlds in which such interventions happen \emph{sufficiency worlds}. A relevant notion present in the literature is that of the probability of sufficiency, which mirrors that of necessity:
\begin{align*} \mathit{PS}(C=c^\star, Y=y^\star) = P(Y_{c^\star} = y^\star \vert C=c', Y=y')\end{align*}
Coming back to Alice, suppose she isn't female, and her loan was approved. What's the probability of her application being rejected if she were female? Again, we need to generalize to keep the other feature fixed here, effectively asking: (1) imagine Alice's score was still insufficient, but she was male and her application succeeded. What's the probability that if she were female, she wouldn't succeed? (2) Imagine she were still a female, but her score was sufficient and she succeeded. What's the probability that she'd have succeeded had her score been insufficient? In the first scenario, her credit score is still insufficient, so she was just one of the lucky 5\% of male applicants whose credit score has not been evaluated. If she were female, the probability of her getting the loan would therefore be 0. So the probability of sufficiency for her gender is 1.

This captures the intuition that her gender played a strong role. Alternatively, suppose her credit score is sufficient, she is still female, and she did get the loan. This means she's among the lucky 90\% of the lucky 20\% women whose credit score has been evaluated. The probability of her application being rejected if her credit score was insufficient is 1 --- which, again, aligns with our intuitions that an insufficient credit score is, in the case of a female applicant, completely sufficient for the application failure.\footnote{How about the probabilities of sufficiency for Ben? If he were female and got the loan, the only way this could have happened is by having a sufficient credit score (and the probability of this outcome is 0.18, i.e. the probability of rejection is  .82). Then, his being male makes the probability of rejection drop to .1, which suggests the probability of sufficiency for him being male is low, as expected. Now suppose his credit score is sufficient, and he got the loan. What's the probability of his application being rejected if we intervene to make his credit score insufficient?  That's 1- (.9 * .05) = .955. So the probability of sufficiency for her credit score is fairly high.}
\begin{align*}
\mathit{PS}^{\text{alice}}(\mathtt{acc} = 0, \mathtt{gen} = 1) 
  &= P\!\left(
      \mathtt{acc}_{\mathtt{gen} = 1} \neq 1 \;\middle|\;
      \mathtt{gen} = 0,\, \mathtt{cred} = 0
    \right) = 1 \\
\mathit{PS}^{\text{alice}}(\mathtt{acc} = 0, \mathtt{cred} = 0) 
  &= P\!\left(
      \mathtt{acc}_{\mathtt{cred} = 0} \;\middle|\;
      \mathtt{gen} = 1,\, \mathtt{cred} = 1
    \right) = 1
\end{align*}


\vspace{2mm}
\noindent \textbf{Probability of necessity and sufficiency.} In principle, we could experiment further, moving on to Pearl's notion of the probability of necessity and sufficiency ($\mathit{PNS}$) that is somehow meant to gauge the two factors simultaneously:
\begin{align*}
PNS  & =  P(y_c, y'_{c'})
\end{align*}
\noindent The disadvantage of doing so, however, would be that this notion does not direct us as to how to build in contextual information to allow for local responsibility attribution. We already had to navigate around this in our toy case above by conditioning on additional information.  Clearly, simple conditioning here wouldn't work in general, as it would bias the outcome one way or another. For instance, suppose $C, Y$ are independent in a simple two-node model with no stochasticity. If you condition on $C =c, Y = y$, then $Y_{c'}$ will always remain $y$. If we condition on $C=c', Y=y'$, then $Y_{c}$ will always remain $y'$, and it is at this point not quite clear how to properly build both upstream context and the factual outcome value correctly into local responsibility attribution. Moreover, conditioning on all elements of the context other than the one we intervene on nullifies the idea that we somehow should perform a search across possible contexts.


\vspace{2mm}
\noindent \textbf{The tale of two worlds.} One lesson from the above is that a properly nuanced picture of the causal role of a feature in the context of a binary outcome should be two-sided, capturing how changing the feature could change the outcome, but also gauging the extent to which having the feature as it is suffices for the outcome.


So, there are some further tasks ahead of us: we need to make sure that we can break the symmetries we already discussed by incorporating search through witnesses in a way analogous to the one used in actual causality, and we need to allow for continuous variables, both as inputs and as outcomes. Moreover, we need to ensure that the resulting notion of responsibility is computationally feasible as models become realistically large. All of this while being Bayesian and probabilistic about our causal impact estimates. This is what we will be moving towards. First, let's focus on probabilistic responsibility attribution analogous to actual causality, which is like PNS in paying attention to both necessity and sufficiency worlds and being probabilities, but like actual causality in paying attention to witnesses and allowing for the symmetry-breaking we've already illustrated.







\section{Probability of Causation: Definitions} \label{sec:definitions}



The above description is quite informal, and we make it more concrete by moving forward within the framework of Structural Causal Models:

\begin{definition}[Causal Model]
    A causal model is a triple \(M = \langle \mathbf{U}, \mathbf{V}, \mathbf{F} \rangle\) where:
\begin{itemize}
    \item $\mathbf{U}$ is a set of background (i.e., exogenous) variables  that are determined by factors outside the model,
    \item $\mathbf{V}$ is a set \(\{V_1, V_2, \ldots V_n\}\) of  variables whose values are determined by other variables already in the model (these are called endogenous variables),
    \item $\mathbf{F}$ is a set of functions \(\{f_1, f_2, \ldots, f_n\}\) such that each \(f_i\) is a mapping from \(\mathbf{U_i} \cup \mathbf{\mathit{Pa}_i}\) to \(\mathbf{V_i}\) where \(\mathbf{U_i} \subseteq \mathbf{U}\), \(\mathbf{\mathit{Pa}_i} \subseteq \mathbf{V} \backslash \mathbf{V_i}\) and the entire set forms a mapping from $\mathbf{U}$ to $\mathbf{V}$. In other words, for each $v_i$, its associated function  \(f_i\) 
    assigns a value to \(\mathbf{V_i}\) that depends on an appropriate set of "upstream" variables in \(\mathbf{V} \cup \mathbf{U}\): \[v_i = f_i(\mathbf{\mathit{Pa}_i}, u_i), \;\;\; i = 1, 2, \ldots n\]
    so that the entire set \(\mathbf{F}\) has a unique solution.
\end{itemize}
\end{definition}

\noindent Such models can be  equipped with probabilities to define probabilistic causal models as follows~\cite{causalityPearl}:
\begin{definition}[Probabilistic Causal Model]
    A probabilistic causal model is a pair \(\langle M, P(\mathbf{u}) \rangle\), where \(M\) is a causal model and \(P(\mathbf{u})\) is a probability function defined over the domain of \(\mathbf{U}\).
\end{definition}

This definition allows us to associate probability distributions with causal models. That is, given a distribution $P(\mathbf{u})$ over the exogenous nodes, we have a pushforward joint distribution \(P_j\) on all the model variables \(\mathbf{U} \cup \mathbf{V}\) via the chain rule, noting that $P(\mathbf{v} \mid  \mathbf{u})$ is the indicator $\mathbbm{1}\left ( (M, \mathbf{u}) \vDash \mathbf{v}\right)$:
\begin{align} \label{eq:joint} P_j(\mathbf{u}, \mathbf{v}) = P(\mathbf{u}) \cdot \mathbbm{1}\left ( (M, \mathbf{u}) \vDash \mathbf{v}\right)\end{align}
where \(\mathbf{u}\) and \(\mathbf{v}\) are assignments to \(\mathbf{U}\) and \(\mathbf{V}\) respectively. We can also recover marginal probabilities for any model variables $\mathbf{a} \subseteq \mathbf{U}\cup \mathbf{V}$ as follows
\[
P(\mathbf{a}) = \sum_{\mathbf{b}} P_j(\mathbf{a}, \mathbf{b}), \mbox{  where } \mathbf{b} = \mathbf{U}\cup \mathbf{V} \setminus \mathbf{a}.
\]


\begin{definition}[Causal set]
    Given a probabilistic causal model, \linebreak  \(\langle \mathbf{U}, \mathbf{V}, \mathbf{F}, P(\mathbf{u})\rangle\), a causal set is a set of variable value pairs \(\{\mathbf{X} = \mathbf{x} ~|~ \mathbf{X} \subseteq  \mathbf{V}, \mathbf{x} \in \mathit{dom}(\mathbf{X})\}\). 
\end{definition}
\noindent As such, a causal set is just a set of variables and some values; what makes it causal is that such sets will be considered potential causes or will be intervened on. 
We define interventions in causal models in alignment with Pearl's do calculus~\cite{causalityPearl}  as follows:

\begin{definition}[Interventions]
    Given a causal model \(\langle \mathbf{U}, \mathbf{V}, \mathbf{F} \rangle\), variable \(X \in \mathbf{V}\) can be intervened with value \(x \in \textit{dom}(X)\) by replacing the function for \(X\) by \(X = x\). We use the notation \([X \leftarrow x]\) to denote the intervention.
\end{definition}

\noindent A short notational reminder before we move forward. In what follows, causal sets will take on multiple roles. At times, they will represent sets of \emph{potential} causes---whose subsets may be intervened upon---in which case we will call them suspects, and denote the set of all such potential causes whose causal role one desires to evaluate by $\mathbf{S}$. When referring to specific sets we \emph{in fact} intervene on, we use $\mathbf{C}$, calling them \emph{active witnesses}.

Analogously, some variables will serve as \emph{potential} elements of the context to be kept fixed at the actual value, denoted by $\mathbf{W}$, and called \emph{witnesses};  when referring to context variables that \emph{in fact} are held fixed in a particular case, we use $\mathbf{T}$ and call them \emph{active witnesses}.


The definitions of sufficiency and necessity below are similar to Pearl's with two main differences. We do not consider the values of the concerned variables in the factual world. This is mainly because we intend to use these definitions in probabilistic causal models across different possible factual worlds, while it is easy to test factivity separately if need be. We also add context sets,  similar to contingencies in Halpern's definition of actual causality~\cite{actualCausalityHalpern}, to incorporate the intuition that token-level causality should be context-sensitive.




\begin{definition}[Sufficiency]
    Given a causal model \(M\) and an assignment \(\mathbf{u}\) to exogenous variables, the causal set \(\{\mathbf{C} = \mathbf{c}\}\) with the context set $\mathbf{T}$ is \textit{sufficient} for outcome \(\varphi\) if the following holds:
    \begin{align}
    \tag{$S_{\mathbf{C} = \mathbf{c}, \mathbf{T}}(\mathbf{u}, \varphi)$}
    (M, \mathbf{u}) 
    \vDash [\mathbf{C} \leftarrow \mathbf{c}, \mathbf{T} \leftarrow \mathbf{t}^\star]\varphi
    \end{align}
     where $\mathbf{t}^\star$ are the actual values of $\mathbf{T}$, that is, \((M, \mathbf{u}) \vDash \{\mathbf{T} = \mathbf{t}^\star\}\). Assuming the model $M$ is held fixed,  we refer to this condition as \(S_{\mathbf{C} = \mathbf{c}, \mathbf{T}}(\mathbf{u}, \varphi)\).

\end{definition}

\begin{definition}[Necessity]
    Given a causal model \(M\) and an assignment \(\mathbf{u}\) to exogenous variables, the causal set \(\{\mathbf{C} = \mathbf{c}\}\) with the context set $\mathbf{T}$ is \textit{necessary} for outcome \(\varphi\) if the following holds for some  \(\mathbf{c'} \in \mathit{dom}(\mathbf{C})\setminus \mathbf{c}\) :
    \begin{align} \label{eq:nec}
        \tag{\(N_{\mathbf{C} = \mathbf{c}, \mathbf{T}}(\mathbf{u}, \varphi)\)}
    (M, \mathbf{u}) \vDash [\mathbf{C} \leftarrow \mathbf{c'}, \mathbf{T} \leftarrow \mathbf{t}^\star]\neg \varphi\end{align}
    where \((M, \mathbf{u}) \vDash \{\mathbf{T} = \mathbf{t}^\star\}\).
    %, \mathbf{x'} \neq \mathbf{x}\).
    Holding the model fixed, we refer to the necessity condition as \(N_{\mathbf{C} = \mathbf{t}, \mathbf{T}}(\mathbf{u}, \varphi)\).
\end{definition}

 Halpern's definition uses an existential quantifier in a few locations: a causal attribution is supposed to hold if \emph{there exists} an alternative value of the active suspect and \emph{there exists} a witness set such that an appropriate necessity condition holds.
 Instead, we will soon explicitly use probabilistic distributions and marginalization.


Given the above definitions of sufficiency and necessity, we define conditioning a probabilistic causal model \(\langle \mathbf{U}, \mathbf{V}, \mathbf{F}, P(u) \rangle\) on sufficiency and necessity conditions being satisfied.

\begin{definition}[Cause-conditioned distribution]\label{def:cause-conditioned}
    Given a probabilistic causal model \(\langle \mathbf{U}, \mathbf{V}, \mathbf{F}, P(\mathbf{u}) \rangle\), a causal set \(\{\mathbf{C} = \mathbf{c}\}\), and a set of context nodes \(\mathbf{T}\) to be kept fixed, a cause-conditioned distribution for the outcome \(\varphi\) is:
    \[P^C_{\mathbf{C} = \mathbf{c}, \mathbf{T}, \varphi}(\mathbf{u}, \mathbf{v}) = P_j(\mathbf{u}, \mathbf{v}) \cdot \mathbbm{1}\left(N_{\mathbf{C} = \mathbf{c}, \mathbf{T}}(\mathbf{u}, \varphi)\right) \cdot \mathbbm{1}\left(S_{\mathbf{C} = \mathbf{c}, \mathbf{T}}(\mathbf{u}, \varphi)\right)\]
\end{definition}


\noindent The intuition here is simple: we take the original distribution $P_j(\mathbf{u}, \mathbf{v})$ and we use the corresponding indicators to ``reject'' cases where one of the conditions fails.  The cause-conditioned distribution is unnormalized (but can be normalized if needed).


To compute a cause-conditioned distribution for an outcome, we need a causal set and a set of context nodes. However, when interested in a causal explanation, one often does not know the exact set of causes or wants to explore various levels of context sensitivity and different contexts ---to achieve this, one needs to search through those. This can be achieved by stochastically choosing the causal set and context nodes. To carry out this stochastic choice, one needs to define a distribution over causal sets and context nodes. Moreover, the probabilities involved in this stochastic search will be ``knobs'' regulating the extent to which we prefer smaller or larger sets of causes, and to which we want to be context sensitive by keeping parts of the actual setting fixed.
One plausible idea for a principled way to search through suspects and witnesses is to set some probabilities of each node being activated and use the same suspect activation probability for all suspects independently, and the same witness activation probability independently for the witnesses. Later on, we will argue that this method leads to issues when scaling to sufficiently large models with continuous outcome variables, and we will introduce a different sampling schema, more suitable for those contexts. For now, let us work with the former explication.


\begin{definition}[Suspect auxiliary and bias]
    Given a variable \(C\) in the potential causal set \(\{\mathbf{S} = \mathbf{s}\}\), let an auxiliary binary variable \(I_C\) denote whether \(C\) is intervened with the value specified in the causal set or not. Given the suspect bias \(b_s\), where $-0.5 < b_s < 0.5$ the probability that \(I_C =1\) is \(0.5 - b_s\) (that is, $b_s$ is a bias towards preemption of that intervention).
    
\end{definition}

\begin{definition}[Witness auxiliary and bias]
    Given a variable \(T\) in the potential context set \(\mathbf{W}\), let an auxiliary binary variable \(J_T\) denote whether \(T\) is kept fixed (preempted) at the factual value or not. Given the witness  bias \(b_w\), the probability that \(J_Y=1\)  is \(0.5 + b_w\) (that is, $b_w$ is a bias towards preemption of $T$ with the actual value).
\end{definition}

\noindent For an indicator variable $\iota\in \{I, J\}$ and a set of variables \(\mathbf{B}\), we abbreviate: \(\iota_B = \bigwedge_{A \in \mathbf{B}} \left(\iota _A = 1 \right)\) as the indicator that all $A \in \mathbf{B}$ were intervened on,   and \(\iota_B^- = \bigwedge_{A\in B} \left(\iota_A = 0\right)\) the indicator that no $A \in \mathbf{B}$ were intervened on.

\begin{definition}[Interventional Setting] \label{def:interventional_setting} A causal setting is determined by specifying a subset $\mathbf{C}$ of the nodes occurring in a set of potential causes $\mathbf{S}$, and
    a subset $\mathbf{T}$ of the set of potential context nodes $\mathbf{W}$.
\end{definition}


The subsets $\mathbf{C}$ and $\mathbf{T}$ will be intervened on: $\mathbf{C}$ to the values specified in a given causal set (alternative values in the necessity world, factual values in the sufficiency world), and $\mathbf{T}$ to factual values. Assuming we already have the alternative values for the elements of $\mathbf{S}$, we can thus equivalently specify the variables with the indicators, so that $\iota_{\mathbf{C}} = \iota_{\mathbf{T}} = 1$, and $\iota^{-}_{\mathbf{W} \setminus \mathbf{C}} = \iota^{-}_{\mathbf{W} \setminus \mathbf{T}} = 1$.

Given a suspect bias $b_s$ and witness bias $b_w$, we can further derive a probability distribution over interventional settings, i.e., determine the probability that (exactly) $\mathbf{C}$ is intervened on and (exactly) $\mathbf{T}$ preempted. This is uniquely specified, after deciding to ``break ties'' with preference to suspects over witnesses, as follows:

\begin{definition}[Probability of Interventional Setting]
    Let $p_s = 0.5 - b_s$ be the probability that a suspect node is not preempted, i.e., that a proposed intervention is executed. Let $p_w = 0.5 + b_w$ be the probability that a witness node is preempted. We obtain interventional settings by first selecting $\mathbf{C}$, then selecting $\mathbf{T}$ from $\mathbf{V} \setminus \mathbf{C}$. Therefore, indicators are chosen as 
    \[
    \begin{array}{ccc}
    \begin{array}{l}
    I_i \sim \mathrm{Ber}(p_s) \\
    \tilde{J}_i \sim \mathrm{Ber}(p_w) \\
    J_i := \tilde{J}_i (1 - I_i)
    \end{array}
    & \Longleftrightarrow &
    \begin{tabular}{lcc}
    \hline
     & $J_i = 0$ & $J_i = 1$ \\
    \hline
    $I_i = 0$ & $(1 - p_s)(1 - p_w)$ & $(1 - p_s)p_w$ \\
    $I_i = 1$ & $p_s$ & $0$ \\
    \hline
    \end{tabular}
    \end{array}
    \]

    \noindent Thus, the probability of any particular interventional setting is:
\begin{align}\label{eq:psetawct}
P^{\mathit{set}}_{b_s,b_w, \{\mathbf{S}=\mathbf{s}\}, \mathbf{W}}(C,T) & = P_{b_s, b_w}\left(I_{\mathbf{C}} \wedge I^-_{\mathbf{S} \backslash \mathbf{C}}  \wedge  J_\mathbf{T}  \wedge J^-_{\mathbf{W} \backslash \mathbf{T}} \right)   % \left( (0.5 - a)^{|C|} (0.5 + a)^{|\mathbf{X}\backslash C|} \right) \times \left( (0.5 + w)^{|T|}(0.5 - w)^{|\mathbf{W}\backslash T|} \right)
\end{align}
\end{definition}
\noindent 



        
    % \begin{flalign*}
    % P_{b_s, b_w}\bigg(\bigwedge_{c \in C} I_c \wedge \bigwedge_{c \in \mathbf{X} \backslash C} &\neg I_c \bigwedge_{t \in T} J_t \wedge \bigwedge_{t \in \mathbf{W} \backslash T} \neg J_t\bigg) = \\ & (0.5 + a)^{|C|} (0.5 - a)^{|\mathbf{X}\backslash C|} \cdot (0.5 + w)^{|T|}(0.5 - a)^{|\mathbf{W}\backslash T|}
    % \end{flalign*}
% \end{definition}


Next, we want to be able to weigh cause-conditioned distributions by the interventional setting probabilities. The result is called an explanation-extended distribution.

\begin{definition}[Explanation-extended distribution]
    Given a probabilistic causal model \(\langle \langle \mathbf{U}, \mathbf{V}, \mathbf{F}\rangle, P(\mathbf{u})\rangle\), an outcome \(\varphi\), a set of possible causes \(\{\mathbf{S} = \mathbf{s}\}\) and its subset \(\mathbf{C}\), a set of context nodes \(\mathbf{W}\) and its subset \(\mathbf{T}\), suspect bias \(b_s\) and witness bias \(b_w\), the explanation-extended distribution can be defined as follows:

    \begin{flalign*}
    P^E_{b_s, b_w, \varphi, \{\mathbf{S} = \mathbf{s}\}, \mathbf{W}}(\mathbf{u}, \mathbf{v}, \mathbf{C}, \mathbf{T}) = 
    &P^{\mathit{set}}_{b_s,b_w, \{\mathbf{S}=\mathbf{s}\}, \mathbf{W}}(\mathbf{C},\mathbf{T})  P^C_{\mathbf{C}, \mathbf{T}, \varphi} (\mathbf{u}, \mathbf{v})
    \end{flalign*}
\end{definition}


\noindent This means that for any particular interventional setting $(\mathbf{C},\mathbf{T})$, we weigh its cause-conditioned distribution by the setting probability distribution corresponding to $(\mathbf{C},\mathbf{T})$.
Note again that the explanation-extended distribution as defined above is unnormalized.

At this point, we have the main components of our approach available. To recap: we start with a joint probability in the baseline model. We multiply it by the necessity and sufficiency indicators to obtain a \emph{cause-conditioned} (unnormalized) distribution. Then we weight it by an interventional setting probability to obtain an \emph{explanation-extended} (unnormalized) distribution. The latter can be used to explicate a range of quantities of interest, to which we will now move. 

The previous subsection defined the notion of an explanation-extended distribution, which can be used to provide explanations for an outcome \(\varphi\). We first define the probability of causation and the degree of responsibility, and elaborate on how these constitute explanations in this setup.

\begin{definition}[Probability of Causation]
    The probability that a causal set \(\mathbf{C}\) causes an outcome \(\varphi\) in a probabilistic causal model \(\langle \mathbf{U}, \mathbf{V}, \mathbf{F}, P(\mathbf{u})\rangle\) given the set of potential causes \(\{\mathbf{S} = \mathbf{s}\}\), a set of potential context nodes \(\mathbf{W}\),  and suspect and witness biases \(b_s\) and \(b_w\) is defined as follows:
    \[P(C \rightarrow \varphi) = \sum_{\mathbf{u}, \mathbf{v}} \sum_{\mathbf{T} \subseteq \mathbf{W}}P^E_{b_s, b_w, \varphi, \{\mathbf{S} = \mathbf{s}\}, \mathbf{W}}(\mathbf{u}, \mathbf{v}, \mathbf{C}, \mathbf{T}) \]
\end{definition}


\noindent That is, to evaluate the ``causal power'' of a set $\mathbf{C}$  over an outcome $\varphi$, we (i) intervene with $C$, (ii) marginalize all variables, (iii) marginalize the choice of which context set is kept fixed, and estimate the resulting explanation-extended probability of $\varphi$.


% All the above definitions are compatible with causal models with continuous distributions as well. But these definitions can become infeasible when the outcome \(\varphi\) is a measure-zero event where a continuous random variable exactly takes on a particular value.

% This particularly messes with Definition~\ref{def:cause-conditioned} where we create a distribution conditioned on sufficiency and necessity, two measure-zero events. To counter this, we propose that instead of using indicator variables for conditioning, we score them using soft equality and non-equality functions. We first provide the definitions of these score functions and then we provide the extended version of a cause-conditioned distribution.

% \begin{definition}[Equality functions]
%     An equality function \(f_= : \mathbb{R} \times \mathbb{R} \rightarrow \mathbb{R}\) is a monotonically increasing function of \(d = |r_1 - r_2|\) where \(r_1, r_2\) are its two inputs. 
% \end{definition}

% \noindent An example of an equality function would be \[f_=(r_1, r_2) = \log{(\mathit{pdf}(|r_1 - r_2|, N(0, \sigma)))}\] where \(\mathit{pdf}(r, N(0, \sigma))\) denotes the pdf value of a zero-centered normal distribution at point \(r\).
% {\color{red}[DW: Hmm? This is monotonically \emph{decreasing} in $d$...]}

% \begin{definition}[Inequality functions]
%     An inequality function \(f_{\neq} : \mathbb{R} \times \mathbb{R} \rightarrow \mathbb{R}\) is a monotonically decreasing function of \(d =|r_1 - r_2|\) where \(r_1, r_2\) are its two inputs. 
% \end{definition}

% \noindent An example of an inequality function would be \[f_{\neq}(r_1, r_2) = \log{(2 \cdot \mathit{cdf}(|r_1 - r_2|, N(0, 0.1)) - 1)}\] where \(\mathit{cdf}(r, N(\mu, \sigma))\) denotes the cdf value of a normal distribution at point \(r\).


% In our implementation we go a bit further, as we parametrize the functions with a scaling factor $a$ and we want scores to be assigned in the log probability space. So we further require the following:
% \begin{itemize}
% \item  As $a$ tends to zero, $f_{eq}^a$ tends to the value it takes when $r_1=r_2$ (call it $\tau$, only in that case, and
% to negative infinity otherwise, 
% \item As $a$ goes to infinity $f_{eq}^a$ tends to $\tau$.

% \item As $a$ tends to allowed minimum (which might be finite, depending on  the functional form of $f_{neq}$), the difference between outcomes $f_{neq}^a(r_1,r_1)$ and $f_{neq}^a(r_1,r_2)$ where $r-1 \neq r_2$ goes to negative infinity.

% \item As $a$ goes to infinity, $f_{neq}^a(r_1,r_2)$ tends to that of $f_{neq}^a(r_1,r_1)$.

% \item For any scale $a$, $f_{neq}^a(r_1,r_2)$ tends to zero
%      $d$ tends to infinity, to the minimum of its range as $d$ tends to zero, and doesn't equal the minimum if $d\neq 0$.
% \end{itemize}



% \begin{definition}[Cause-Conditioned Distribution, Extended]
%     Given a probabilistic causal model \(\langle U, V, F, P(u) \rangle\), a causal set \(\{\mathbf{X} = \mathbf{x}\}\), and a set of context nodes \(\mathbf{W}\) to be kept fixed, a cause-conditioned distribution for the outcome \(\varphi := (R = r)\) where \(R \in V, r \in \mathit{dom}(R)\) and \(R\) is continuous, if the following hold:
%     \[(M, \mathbf{u}) 
%     \vDash [\mathbf{X} \leftarrow \mathbf{x}, \mathbf{W} \leftarrow \mathbf{w}^\star](R = r_s)\]
%     \[(M, \mathbf{u}) 
%     \vDash [\mathbf{X} \leftarrow \mathbf{x'}, \mathbf{W} \leftarrow \mathbf{w}^\star](R = r_n)\]
%     then cause-conditioned distribution can be defined as follows where \(f_=\) and \(f_{\neq}\) are equality and inequality functions respectively:
%     \[P^C_{\mathbf{X} = \mathbf{x}, \mathbf{W}, \varphi}(\mathbf{u}, \mathbf{v}) = P_j(\mathbf{u}, \mathbf{v}) \cdot f_{\neq}(r, r_n) \cdot f_{=}(r, r_s)\]
% \end{definition}
% {\color{red}[DW: To me, the somewhat simpler approach would be to simply be a smooth analogue to the indicator function, e.g., $\exp(-r/c)$ for large $c$. What are the implications of using this over the equality/inequality functions? Theoretically, it avoids negative ``probabilities,'' which always weirds people out :)]}

% \noindent The remaining  definitions remain the same, with only the extended definition of cause-conditioned distribution swapped for the original one.

% Taking a step back,  the probabilities of causation, while shaped by the original model itself, are also impacted by what the user does with a few knobs that they might manipulate. Manipulating the probabilities of intervention preemption mirrors user preference for (or against) smaller cause sets. Manipulating the probabilities of witness preemption mirrors how context-sensitive the evaluation is supposed to be. Adding parameterized equality and inequality functions allows the user to decide how steeply the distances between outcomes should be rewarded or penalized by the respective equality/non-equality functions. 


\section{Relation with Actual Causality} \label{sec:relation_with_ac}

The notion of explanation-extended distribution is not meant to explicate Halpern's notion of actual causality; rather, the goal was to combine the intuitions underlying Halpern's notion with those motivating Pearl's notion of probability of necessity and sufficiency to provide a more general tool for causal attribution. Nevertheless, to show that we preserve the spirit of Halpern's definition, in this section, we discuss the connection between actual causality and the notions of explanation-extended distribution (for proofs, see the appendix). The goal is to show that there is a natural choice of parameters for our method that makes it give verdicts that agree with actual causality. With one caveat: probability of causation maximization will not identify subset-minimal causes that are not cardinality-minimal.



Let's start with the main definition used by Halpern:

\begin{definition} \label{def:ac}
    \(\mathbf{C} = \mathbf{c}\) is an actual cause of \(\varphi\) in the causal setting \((M, \mathbf{u})\) if the following three conditions hold:
    \begin{itemize}
        \item \textbf{Factivity}: \((M, \mathbf{u}) \vDash (\mathbf{C} = \mathbf{c}\,)\) and \((M, \mathbf{u}) \vDash \varphi\),
        \item \textbf{Context-sensitive Necessity}: There is a set of variables \(\mathbf{T} \subseteq \mathbf{V}\) and an alternative setting \(\mathbf{c'}\) of variables in \(\mathbf{C}\) such that if \((M, \mathbf{u}) \vDash \mathbf{T} = \mathbf{t^\star}\), then
            \[(M, \mathbf{u}) \vDash [\mathbf{C} \leftarrow \mathbf{c'}, \mathbf{T} \leftarrow \mathbf{t^\star}] \neg \varphi,\]
        \item \textbf{Minimality}: \(\mathbf{C}\) is minimal; there is no strict subset \(\mathbf{C_{sub}} \subset \mathbf{C}\) such that \(\mathbf{C_{sub}} = \mathbf{c_{sub}}\) satisfies the above two conditions where $\mathbf{c_{sub}}$ is the restriction of \(\mathbf{c}\) to the variables in \(\mathbf{C_{sub}}\). 
    \end{itemize}
\end{definition}


First, we notice that if the actual causality holds, there are some fairly natural settings of suspects and witnesses such that for any (regular) witness and suspect biases, the resulting integrated probability of causality (still parameterized by noise) is greater than zero.


\begin{theorem}\label{th:ac-exp}
    Suppose that $\mathbf{C}$ is an actual cause with a witness set $\mathbf{T}$ in $(M, \mathbf{u})$.  Take any  $\{\mathbf{S} = \mathbf{s}\} \supseteq \mathbf{C}$ as the potential cause set, and follow the definition in taking the set of all endogenous variables $\mathbf{V} \supseteq \mathbf{T}$ as the potential witness set. Then for any suspect bias  $ - 0.5 < b_s < 0.5$, and  any witness  bias $ - 0.5 < b_w < 0.5$ we have: \begin{align}
    P^{E}_{b_s, b_w , \varphi, \{\mathbf{S} = \mathbf{s}\}, \mathbf{V}}(\mathbf{u}, \mathbf{v}, C, T) > 0 \label{eq:pect}\\
    \text{ and therefore } \nonumber\\ 
\sum_{\mathbf{T} \subseteq \mathbf{V}}\, P^E_{b_s, b_w , \varphi, \{\mathbf{S} = \mathbf{s}\}, \mathbf{V}}(\mathbf{u}, \mathbf{v}, C, T) > 0 \end{align}
\end{theorem}


There is also a clear connection in the opposite direction: if the same noise-parameterized probability of causation is greater than zero, the necessity condition in the definition of actual causality holds.


\begin{theorem}\label{th:exp-ac} Moreover, if  
$\sum_{T \subseteq \mathbf{V}} P^{E}_{b_s, b_w, \varphi, \mathbf{C}, \mathbf{V}}(\mathbf{u}, \mathbf{v}, \mathbf{C}, \mathbf{T}) > 0$,  the context-sensitive necessity condition in Definition \ref{def:ac} holds. \label{th:if_pr_then_nec}
\end{theorem}






With minimality, the situation is somewhat more complicated. For one thing, our choice of suspect and witness preemption biases will have an impact on probability-preferred set sizes. On the other hand, the way that  $P^{\mathit{set}}_{b_s,b_w, \{\mathbf{S}=\mathbf{s}\}, \mathbf{V}}(\mathbf{C}, \mathbf{T})$ is defined makes sets comparable in light of their cardinalities, whereas the minimality Definition \ref{def:ac}  is formulated in terms of the subset relation, which is bound to produce some misalignment. Some useful observations can be made, though.

One of them is that there is a somewhat degenerate simplified setting in which positive explanation-extended probability, on the assumption of factivity, entails minimality wrt. to the same parameters: it is enough to assume null witness bias and that the suspect bias prefers smaller explanations. The strategy used in the proof entails that we can safely marginalize out the witnesses as well.


\begin{theorem} \label{th:local_max_to_ac} Write $\Phi (\mathbf{C}, \mathbf{T})$ for $P^{E}_{b_s, b_w, \varphi, \mathbf{C}, \mathbf{V}}(\mathbf{u}, \mathbf{v}, \mathbf{C}, \mathbf{T})$. Consider the set of optimizers  $\Gamma = \arg\max_{\mathbf{C} \subseteq \{\mathbf{S} = \mathbf{s}\}, \mathbf{T} \subseteq \mathbf{V}} \Phi(\mathbf{C}, \mathbf{T})$.  Assume $(\mathbf{C^\dagger}, \mathbf{T^\dagger})\in \Gamma$ with $\Phi(\mathbf{C^\dagger}, \mathbf{T^\dagger})>0$ and assume further that factivity is satisfied for $\mathbf{C^\dagger}$ (as in Definition \ref{def:ac}). Let the witness bias be $b_w=0$ and let the suspect preemption bias be $b_s>0$ (that is, we positively prefer suspect preemption, thus building preference for smaller cause sets into our search). Then $C^\dagger$ is an actual cause of $\varphi$ according to Definition \ref{def:ac}.
\end{theorem}



\begin{corollary}\label{cor:marginalize} Let's keep the notation from Theorem \ref{th:local_max_to_ac}. Consider the set   of optimizers  $\Delta = \arg\max_{\mathbf{C} \subseteq \{\mathbf{S} = \mathbf{s}\}} \sum_{\mathbf{T} \subseteq \mathbf{V}} \Phi(\mathbf{C}, \mathbf{T})$.  Assume $\mathbf{C^\dagger}\in \Delta$ with \[\sum_{\mathbf{T} \subseteq \mathbf{V}} \Phi(\mathbf{C^\dagger}, \mathbf{T})>0\] and assume still that factivity is satisfied for $\mathbf{C^\dagger}$. Keep $b_w=0$ and $p_s>0$  as in Theorem \ref{th:local_max_to_ac}. Then $\mathbf{C^\dagger}$ is an actual cause of $\varphi$ according to Definition \ref{def:ac}.
\end{corollary}



Now we can return to the misalignment we have mentioned and elaborate on the nuance. One can imagine a reasoning mode when one looks for causes of an outcome by inspecting probabilities of causation $P(C \rightarrow \varphi)$ not knowing what the values of $\mathbf{u}$ are, but knowing that $\varphi$ occurred. It is worth pointing out that  $\arg\max_{\mathbf{C} \subseteq \mathbf{X}}P(\mathbf{C} \rightarrow \varphi)$ might fail to include some $\mathbf{C}$ such that for some setting $\mathbf{u}$, $\mathbf{C}$ is nevertheless an actual cause of $\varphi$.


\begin{observation}\label{obs:misalignment} Define $\Xi = \arg\max_{\mathbf{C} \subseteq {X}} P(C \rightarrow \varphi)$. It is not the case that if $C^\dagger$ is an actual cause in $(M, \mathbf{u})$ for some $\mathbf{u}$, $C^\dagger \in \Xi$.
\end{observation}

The rather straightforward reason for this is that actual cause sets might be subset-minimal without being cardinality-minimal, while the probability of causality maximization will drop those that are not cardinality-minimal.

This can be relatively easily alleviated by not maximizing the probability of causation if one searches for actual causes, but rather finding all cause sets with non-zero probabilities (by our previous theorems, all of those should satisfy the necessity requirement) and subsequently filtering for subset-minimality.

However, we find that some idiosyncratic components of the notion of actual causality don't make it directly scalable anyway, so at this point, we wrap up a discussion of that notion and move to the discussion of a related, more probabilistic, and more scalable notion.


% \begin{observation} Consider 
% \begin{align*}\Xi & = \left\{C \vert \left(\sum_{\mathbf{T} \subseteq \mathbf{V}} \Phi(C, T)>0\right),  \neg \exists C' \subset C \left( \sum_{\mathbf{T} \subseteq \mathbf{V}} \Phi(C', T)\right) \right\}
% \end{align*}

% \begin{enumerate} 
% \item $\Delta \subseteq \Xi$.
% \item If $C^\dagger \in \Xi$, $C^\dagger$ is an actual cause. 
% \end{enumerate}
% \end{observation}



% \begin{proof}
% TODO
% \end{proof}


% \subsection{Factivity}

% Our definition of explanation-extended distribution does not enforce the criteria of factivity but that can be easily achieved by conditioning. We also believe that the factivity is a restricting condition since we do not wish to limit our causal reasoning to a particular factual world. We wish to generalize to other possible worlds as well.



% \section{Questions}
% \begin{itemize}
%     \item These definitions do not exactly coincide with Halpern and Pearl's definitions. How do we want to go about justifying these definitions?
%     \item Are there any other quantities that people think we can define above? 
% \end{itemize}




\section{Scaling up and continuous variables}
\label{sec:scaling}


\vspace{2mm}
\noindent \textbf{Causal setting sampling.} While there is some plausibility to thinking of witness or suspect candidate sampling independently with equal probability when the model is relatively small, when dealing with a large model, this will have the unintended consequences of strongly biasing the size of the activated sets towards the expected value of the associated binomial distribution. 

For instance, suppose your sets of witness candidates and of suspects are disjoint, each containing 24 variables. Suppose further you want to obtain a sample from the expanded model, which would allow you to perform multiple rejection sampling steps to compute some causal impact score for each of the suspects. That is, for each suspect variable $X$, you would like to be able to condition on $X$ being intervened as a suspect (i.e., its suspect indicator being true). For this reason, you keep the suspect preemption probability not too high, say starting with setting it to .5. As a side effect, the number of suspects intervened simultaneously will center around the expected of the relevant binomial, i.e. around $0.5 \times 24 = 12$, and so the scoring will be mostly based on rather large intervention sets, which is unintended and in realistic settings will make your evaluation noisy. Similarly with the witnesses, now the size of the context set kept fixed will be around 12, which means you will predominantly force the search to drop half of the context, which might again be unintended and lead to a lot of noise. 

So we need to generalize, allowing for other sampling schemata for suspect and witness activation, so that $C \sim \Gamma_s(2^S)$ and $T \sim \Gamma_w(2^W)$, where $\Gamma_s$ and $\Gamma_w$ are suspect and witness sampling distributions, with the dependency excluding variables from being both. We also get the associated probability functions $p^\Gamma_s$, $p^\Gamma_w$. One pattern that is intuitive and worked well in a real-life case we describe in the Appendix, is setting an upper limit on the number of activated antecedents, effectively assuming that we care only about explanation sets or cause interactions of certain complexity, sampling the size uniformly from the allowed range, and then sampling that many variables uniformly from the suspects set (and symmetrically for witness candidates, we put a cap on the number of context elements to be dropped, sample a size uniformly from the remaining options, and sample that many active witnesses uniformly from the witness candidates diminished by the active suspects.

\vspace{2mm}
\noindent \textbf{Selecting alternative values.} Causal setting sampling will choose variables to intervene on. We also need to decide what values to intervene with.
We need to go beyond assuming the variables to be binary and work out a more general semantics for the suspect in `if $X$ had a different value,...'. In general, for any selection $C$ of active antecedents, we need to sample alternative values from some joint alternative value distribution $\Delta_a$. The reason why we need a joint distribution is that if in the search more than one suspects are active, the way they can be different might depend on each other: for instance, an alternative in which an apartment is larger and yet has fewer rooms and fewer bathrooms doesn't seem like a very likely scenario and should be considered less often than the one in which the combination of values is more natural.


One example of a natural candidate here is to assume that in the context of continuous variables, this means `a sufficiently different value' with some threshold $\epsilon$ and taking the sample from the model itself, where the relevant distribution has been excised to exclude the region of $\pm \epsilon$ from the factual. 

\vspace{2mm}
\noindent \textbf{Causal impact estimation.} 
The steps to causal impact estimation on this approach are as follows. First, we sample the active suspects and witnesses:
\begin{align}
\mathbf{C} \sim \Gamma_s(2^S) & & \mathbf{T} \sim \Gamma_w(2^W)
\end{align}
\noindent Next, we sample alternative interventions for $\mathbf{C}$: 
\begin{align}
\mathbf{c'} \sim \Delta_a
\end{align}
Then the sufficiency intervention, $s_{C,T}$ is obtained by setting the values of $\mathbf{C}, \mathbf{T}$ to their factual values, $\mathbf{c}^\star, \mathbf{t}^\star$, while the necessity intervention, $n_{C, T}$, is obtained by setting the values of $\mathbf{C}$ to $\mathbf{c'}$ and keeping the values of $\mathbf{T}$ fixed at $\mathbf{t}^\star$:
\begin{align}
s_{C,T}  & = do(\mathbf{C} = \mathbf{c^\star}, \mathbf{T} = \mathbf{t^\star}) \\
n_{C,T} & = do(\mathbf{C} = \mathbf{c'}, \mathbf{T} = \mathbf{t^\star})
\end{align}
For an outcome of interest $y$, the family $p$ of intervened densities determined by the causal model provides us with two densities under two interventional regimes (where $\mathbf{u}$ is the noise):
\begin{align}
p^s_{C,T}(y\vert  \mathbf{u}) & = p (y \vert \mathbf{u}, s_{C, T})\\
p^n_{C,T}(y\vert \mathbf{u}) & = p(y \vert \mathbf{u}, n_{C,T})
\end{align} The \textbf{expanded intervened densities} are obtained by weighting these by the probabilities of $C, T$ being chosen:
\begin{align}
p^{es}{C,T}(y \vert \mathbf{u}, C, T) & = p^s_{C,T}(y\vert  \mathbf{u})  p^\Gamma_s(C) p^\Gamma_w(T)\\ 
p^{en}(C, T)( y \vert \mathbf{u}, C, T) & = p^n_{C,T}(y\vert \mathbf{u})    p^\Gamma_s(C) p^\Gamma_w(T)
\end{align}

Now, suppose we consider the causal impact of a feature $X_k \in S$. We can condition on its being an active antecedent, giving us a conditioned  \textbf{search density}. Here, the summation completes marginalization across the relevant subsets, as the expanded intervened densities already include multiplication by interventional setting probabilities:
\begin{align}
\overset{?}{p}^s_k(y \vert \mathbf{u})    &= \sum_{X_k \in C \subseteq A, T\subseteq W} p^{es}{C,T}(y \vert \mathbf{u}, C, T) \\ 
\overset{?}{p}^n_k(y \vert \mathbf{u})   & = \sum_{X_k \in C \subseteq A, T\subseteq W}p^{en}(C, T)( y \vert \mathbf{u}, C, T) 
\end{align}
This restricts our attention to cases in which in the necessity world, $X_k$ (together with potentially some other suspect variables)  have been intervened on with an alternative value, and in the sufficiency world, $X_k$ (and the same other causal suspects) have been intervened with the factual value. 


Let $y^s$ denote the outcome in the \emph{sufficiency world} and $y^n$ denote the outcome in the \emph{necessity world}. By construction, these are independent conditional on $\mathbf{u}$, and so we get a joint push-forward distribution:
\begin{align}
\overset{?}{p}(\langle y^s, y^n \rangle \mid \mathbf{u})  & = \overset{?}{p}^s_k(y^s \mid \mathbf{u}) \; \overset{?}{p}^n_k(y^n \mid \mathbf{u}).
\end{align}
\noindent which also generates a push-forward on any causal impact scoring function $ci(y^s, y^n, y^\star)$.  One straightforward function of that sort would equally reward distance from the factual value $y^\star$ in the necessity world, and proximity to the factual value in the sufficiency world. Using the two in the scoring is an echo of Pearl's intuition to use the probability of both sufficiency and necessity.
\begin{align*}
ci(y^s, y^n, y^\star) & = \vert y^n - y^\star \vert   - \vert y^s - y ^\star \vert     
\end{align*}
Assuming the model also produces a distribution over $\vert{u}$, one natural score summary could be the expected value of that function.
\begin{align}\label{eq:score}
\mathbb{E}_{\mathbf{u}}\Big[\, \mathbb{E}_{y^s, y^n \mid \mathbf{u}} \big[ ci(y^s, y^n, y^\star) \big] \,\Big]
= \\  =\int p(\mathbf{u}) 
   \int \int ci(y^s, y^n, y^\star) \; p^s_k(y^s \mid \mathbf{u}) \; p^n_k(y^n \mid \mathbf{u}) 
   \, dy^s \, dy^n \, d\mathbf{u}. \nonumber
\end{align}






\section{Case studies}
\label{sec:case_studies}


The concepts described in this paper have been implemented using ChiRho.  In \ref{sub:synthetic}, we develop a synthetic loop that evaluates the performance of the causal impact scoring method described in Section \ref{sec:scaling} using a generative process that involves overdetermination and undercutting. In \ref{sub:amv}, we apply the same method to a real-life large-scale problem of real estate automated market valuation.



\subsection{Synthetic evaluation with overdetermination and undercutting}
\label{sub:synthetic}


In this evaluation, we build off the existing counterexamples to the but-for clause as an explanation tool. Recall that the  initial intuition one might have is the following:

\begin{quote}
$X=x$ is responsible for $Y=y$ just in case $Y$ would have been different had we intervened on $X$ to be different, i.e. if $Y_{x'} = y'$ for some $x' \neq x, y' \neq y$. 
\end{quote}

A typical counterexample to this is cast in terms of categorical variables \cite{actualCausalityHalpern}: 

\begin{quote} Sally and Billy pick up stones and throw them at a bottle. Sally’s stone gets there first, shattering the bottle. Both throws are perfectly accurate, so Billy’s stone would have shattered the bottle had it not been preempted by Sally’s throw. The intuition is that Sally's throw is responsible for the shattering, but not so Billy's. However, it is false that had Sally not thrown the stone, the bottle would not have shattered, so the simple counterfactual condition fails.
\end{quote}

What we want to investigate is the performance of responsibility attribution, using the same implementation as the one used in Appendix \ref{sub:amv}, following the discussion in Section \ref{sec:scaling}, in the context of a synthetic generalization of this problem. Suppose we use continuous variables in a model with a similar overdetermination and undercutting structure. Is the tool capable of meeting our expectations when it comes to responsibility attribution? 

The answer is positive. We first introduce the model, introduce intuitive qualitative desiderata, and investigate their satisfaction for 500 local explanations, for events generated using the same generative process, with responsibility scores for all features estimated using one sample of size 5000.

The way we generalize is as follows. The input and output variables are continuous. The $A$ variables are like Sally's throw: they have an unblocked impact on the outcome. The $C$ variables are like Billy's throw: if their linear combination is high enough, they are active and can have an impact on the outcome. Otherwise, their impact is blocked. $B$ and $D$ are like `Sally hits' and `Bill hits', except continuous, and $E$ is some final outcome, except now we don't use a disjunction, but rather addition. 


\begin{align*}
\text{For } i = 0, 1, 2: &  \\
&\quad A_i \sim \mathcal{N}(0, 1), \\
&\quad C_i \sim \mathcal{N}(0, 1), \\
\text{and}  & \\
 B & = 0 \cdot A_0 + 5 \cdot A_1 + 10 \cdot A_2, \\
 cs\_dot & = 0 \cdot C_0 + 5 \cdot C_1 + 10 \cdot C_2, \\
 cs\_win &= \mathbf{1}\{ cs\_dot > B \}, \\
 D & = cs\_dot \cdot cs\_win, \\
 E& = B + D.
\end{align*}


\noindent See Figure \ref{fig:synthetic_model} for a visualization of the model structure.


Say we denote the mean causal impact of a variable $V$ by $ci(V)$. We do not expect to see zero estimates for irrelevant properties, as the score also results from sampling other potential causes and witnesses, so there will always be some distance between the outcomes in both sufficiency and necessity worlds. But we do expect the following qualitative/comparative statements to hold:

Most of the time, $ci(A0) < ci(A1) < ci(A2)$, with $ci(A1)$ and $ci(A2)$ closer to each other and potentially the $ci(A1) > ci(A2)$ for cases where $A1 >> A2$.

Similarly, for the $C$ variables, when their impact is not blocked (i.e. if `cs\_win =1`). Otherwise, we expect $ci(C1) \approx ci(C2) \approx ci(C3)$. Given the difference in means between $ci(A0)$ and $ci(A2)$ close to $6$, let's explicate this proximity as the absolute difference between sampled values being $<2$. Mean estimated scores for all the features involved and the frequency of the binary claims satisfaction are as expected, and can be inspected in Figures \ref{fig:synthetic_attribution} and \ref{fig:synthetic_summary}




\begin{figure}[htbp]
\centering
\includegraphics[width= 11cm]{figures/synthetic_causal_model.png}
\caption{Synthetic model involving overdetermination and undercutting used in the evaluation loop.}
\label{fig:synthetic_model}
\end{figure}


\begin{figure}
    \centering
    \includegraphics[width=0.9\linewidth]{figures/synthetic_causal_model_responsibility_attribution.png}
    \caption{Mean attribution scores to all upstream features in the synthetic model, 500 individual explanations, all based on a single search sample of size 5000. When `cs\_win=0', the $C$ variables are not activated, and so we expect total scores for all of them to be similar (red, violet, dark brown on the left-hand side). We still expect the scores for $A0$, $A1$, and $A2$ generally to increase independently of what happens with the $C$ variables, which we see on both sides (blue, light brown, and green on both sides). We also expect the scores for $C$ variables to increase if they are activated (red, violet, dark brown on the right-hand side). These expectations are satisfied.}
    \label{fig:synthetic_attribution}
\end{figure}


\begin{figure}
    \centering
    \includegraphics[width=0.9\linewidth]{figures/synthetic_causal_model_summary.png}
    \caption{Qualitative claims satisfaction frequency for causal attribution in the synthetic model, 500 local explanations all based on a single search sample of size 5000. Most of the time, the ordering of the scores for $A$ variables agrees with the ordering of the coefficients. Most of the time, close proximity of scores for inactive $C$ variables is quite high; most often, the ordering for active $C$ variables agrees with the ordering of the coefficients. We do not, however, want an explanation to track coefficients only, which, for instance, is illustrated by the last statement: if scores are reversed, this is mostly due to the disparate direction of the difference between the variable values.}
    \label{fig:synthetic_summary}
\end{figure}


\subsection{Causally aware probabilistic automated valuation model}
\label{sub:amv}


Our units of analysis were property transactions in the State of New York. The outcome variable $y$ was the transaction amount (in the log-epsilon-standardized space).

We started with a range of a range of categorical inputs $x_{cat} = x_{1}, \dots, x_k$ and continuous inputs in the transformed space, $x_{con} = x_{k+1}, \dots, x_n$ (usually log-epsilon-standardized, standardized-truncated, or standardized, depending on the variable) for continuous variables. 

A causal DAG $\Gamma$ expresses causal assumptions about how $x_1, ...,x_n$ are related in terms of direct causal connections. It was obtained by iterating with experts on one hand and incorporating clear cases of conditional independence test failures, resulting in 99 edges (see Figure \ref{fig:edges}).





In what follows, $\varphi$ with subscripts represents parameters of different models or model components. Each causal kernel, comprised of an input variable $x_i$ with its parents $\mathrm{pa}_i$, has been represented as a neural conditional distribution, so that:
\begin{align}
x_i & \sim \mathrm{Dist}_i(f_i(\mathrm{pa}_i, \varphi_i))
\end{align}
\noindent where $f_i$ is a deep neural network, $f_i(\mathrm{pa}_i, \varphi_i)$ are distribution parameters passed to $\mathrm{Dist}_i$, and $\mathrm{Dist}_i$ is a probabilistic distribution of an appropriate mathematical form (normal, categorical, or truncated normal).\footnote{For root nodes, $\mathrm{pa}_i$ is empty and $f_i$ is not conditioned.} The loss used in optimization was cross-entropy for categorical distributions and average negative log likelihood for the continuous distributions.


The outcome was modeled by a sparse Gaussian Process composed of linear embeddings of categorical variables, so that:
\begin{align}
x_{gp} & = (x_{con}, \eta(x_{cat}, \varphi_{cat}) ) \\
y & \sim \mathcal{GP}\left(\mu(x_{gp}, \varphi_\mu), \kappa(x_{gp}, \varphi_\kappa)\right)
\end{align}
\noindent where $\eta$ is a linear embedding of the categorical variables into the continuous space and both $\mathcal{GP}$ and $\eta$ are optimized jointly.

The resulting $\mathcal{GP}$ is comparable in performance to a CatBoost model trained on the same data (see Figure \ref{fig:eval}).




The structured causal model $\mathfrak{M}$ first runs the causal Markov kernels in a topological order and then passes the resulting $x_i, \ldots x_n$ to the $\mathcal{GP}$, so that the generative model specifies a joint distribution:
\begin{align*}
x_1, \dots, x_n, y &\sim \mathfrak{M}
\end{align*}
\noindent and, more generally, a family of interventional distributions under interventions $\mathbf{I}=\mathbf{i}$, $\mathfrak{M}_{\mathbf{I}=\mathbf{i}}$:
\begin{docode}
    with do($\mathbf{I}=\mathbf{i}$):\\
    \hspace*{2em} $\mathfrak{M}()$
\end{docode}
One use of the model is to facilitate causally aware imputation: conditioning on the available features and resampling missing features in forward runs. These features can be subsequently passed on to downstream components and ultimately to the $\mathcal{GP}$ model.

Another use is to predict the effects of interventions. Consider intervening on a range of variables
$\mathbf{I}= \mathbf{i}$. To estimate the effects of this intervention on a property transaction with features
$\mathbf{X=x}$:
\begin{docode}
with condition$\left( \mathbf{X}= \mathbf{x} \setminus (\mathbf{I}= \mathbf{i^\star} \cup  \mathbf{I_\delta}   = \mathbf{i_\delta^\star}\right):$

\hspace*{2em}  with do($\mathbf{I}=\mathbf{i}$):

  \hspace*{5em} $\mathfrak{M}()$
\end{docode}
\noindent where  $\mathbf{I_\delta} = \mathrm{descendants}(\mathbf{I})$, providing causally aware estimates  (see Figure \ref{fig:counterfactual_sales} for an example).

We also transform the model into the expanded one, in which we evaluate the impact of features from a pre-selected list of potential causal suspects $\mathbf{S}$. In the expanded model, we have:
\begin{docode}
$\mathbf{C} \sim \Gamma_s(2^\mathbf{S})$   sample active antedents \\
$\mathbf{T} \sim \Gamma_w(2^\mathbf{W})$ sample active witnesses \\
with sample\_alternatives(factual = $\mathbf{C} = \mathbf{c^\star}$): \\
\hspace*{2em} $\mathbf{c'} \sim \mathfrak{M}$

\vspace{2mm}

$\mathbf{I}_s   := \{\mathbf{C} = \mathbf{c^\star}\} \cup \{ \mathbf{T} = \mathbf{t^\star}\}$ \\
$\mathbf{I}_n := \{ \mathbf{C} = \mathbf{c^{\prime}} \} \cup \{ \mathbf{T} = \mathbf{t^{\star}} \}$ \\

with MultiWorldCounterfactual: \\
\hspace*{2em}  with do$\left( \langle \mathbf{I}_s, \mathbf{I}_n \rangle \right) $: \\
\hspace*{4em} $\langle y_s, y_n \rangle = \mathfrak{M}$

\vspace{2mm}

scores = $ci(y^s, y^n, y^\star)$

\end{docode}
\noindent where (1) $\Gamma_s(2^\mathbf{S})$ is a search distribution used to sample small subsets of $\mathbf{S}$ (to prefer small causal explanations), 
(2) $\Gamma_w(2^\mathbf{W})$ is a search distribution used to sample large subsets of $W$ (to prefer large elements of the context to remain untouched), (3) \texttt{sample\_alternatives} is a handler that replaces distributions with excised ones which sample around the factual values with renormalized excised densities, (4) \texttt{MultiWorldCounterfactual} allows for running and tracking of multiple possible worlds in one pass forward, and (5) is a scoring function.

This results in a Monte Carlo sample, which we can subset by conditioning on a given suspect $X_i$ being active to obtain a sample of the scores corresponding to that feature, which can be used to estimate the expected value of that score. For our prototype, this estimation converges at around 25k samples, and sampling needs to be performed only once to compute attribution scores for all suspects. See Figure \ref{fig:scores} for an example of attribution scores for a single evaluation, contrasted with the corresponding SHAP values for the same model and the same unit of observation.



\begin{figure}
\begin{center}
\includegraphics[width=11cm]{figures/gp_vs_catboost_eval.png}
\end{center}
\caption{Evaluation of the predictive performance of the downstream GP model trained on New York State data up to 2023, tested on hold-out data from 2023 and 2024, compared to an optimized CatBoost model with the same data split.}
\label{fig:eval}

\end{figure}




\begin{figure}[htbp]
    \centering    \includegraphics[width=10cm]{figures/anonymized_upstream_structure.png}
    \caption{DAG used in the causal model for upstream variables in the causal AVM prototype (variables anonymized due to proprietary considerations). Each causal Markov kernel is modeled using a separate neural distribution.}
    \label{fig:edges}
\end{figure}
    


\begin{figure}[htbp]
    \centering  
    \includegraphics[width=\textwidth]{figures/counterfactual_sale_amounts.png}

    \caption{Counterfactual sale amount predicted by the causal AVM model. Median unintervened prediction (blue), upon the addition of one half-bathroom, moves realistically up (orange) with proper propagation of this change to downstream variables. Furthermore, the flat GP prediction obtained by simply changing the input value for half-bathrooms is overly optimistic about the impact of the extreme intervention of adding 4 half-bathrooms (red). This is in contrast to the negative impact predicted by the causal model (green).}
    \label{fig:counterfactual_sales}
\end{figure}




\begin{figure}[htbp]
\centering
\includegraphics[width = 11cm]{figures/search_query_figure.png}
\caption{Attribution scores for a single evaluation, shown alongside SHAP values from the flat downstream Gaussian Process (mean predictions) for the same unit of observation. The attribution scores estimate the quantity defined in \eqref{eq:score}—roughly, the expected price change in the necessity world, penalized by the change in the sufficiency world, relative to the observed value. In contrast, SHAP places nearly all weight on features that, in the upstream causal model, lie downstream of most other features.}
\label{fig:search}
\end{figure}


\bibliographystyle{plain}
\bibliography{references}




\appendix
\section{Proofs}
\label{app:proofs}


\begin{proof}[Proof of Theorem \ref{th:ac-exp}]  Assuming the antecedent of the theorem, freeze $\mathbf{u}$, $\mathbf{c'}$, and $\mathbf{T}$ as witnesses to the satisfaction of the conditions specified in Definition \ref{def:ac}. For the condition of context-sensitive necessity to be satisfied, $\mathbf{T} \cap \mathbf{C}$ needs to be disjoint, otherwise the intervention with alternative values $\mathbf{c'}$ and at the same time actual values $\mathbf{t^\star}$ would be ill-defined. Further, set $\mathbf{v}$ to be the values that $\mathbf{V}$ take in  $(M,u)$, so that $(M, \mathbf{u}) \models \mathbf{V} = \mathbf{v}$.  
We follow the definition of an explanation-extended distribution to decompose \eqref{eq:pect} as follows:
\begin{align}
    P^{\mathit{set}}_{b_s,b_w, \{\mathbf{S}=\mathbf{s}\}, \mathbf{V}}(C,T)  P^C_{C, T, \varphi} (\mathbf{u}, \mathbf{v}) = \\
      P^{\mathit{set}}_{b_s,b_w, \{\mathbf{S}=\mathbf{s}\}, \mathbf{V}}(C,T) \cdot P_j(\mathbf{u}, \mathbf{v}) \cdot \mathbbm{1}\left(N_{\mathbf{C} = \mathbf{c}, \mathbf{T}}(\mathbf{u}, \varphi)\right) \cdot \mathbbm{1}\left(S_{\mathbf{C} = \mathbf{c}, \mathbf{T}}(\mathbf{u}, \varphi)\right) \label{eq:decomposition}
\end{align}
\noindent The goal now is to show that each of the components is greater than 0.

\begin{itemize}
\item Starting with $P^{\mathit{set}}_{b_s,b_w, \{\mathbf{S}=\mathbf{s}\}, \mathbf{W}}(C,T)$, by \eqref{eq:psetawct}, this is identical to \linebreak  $P_{b_s, b_w}\left(I_c \wedge I^-_{\mathbf{X} \backslash C}  \wedge  J_T  \wedge J^-_{\mathbf{V} \backslash T} \right)$. Given  $\mathbf{T} \cap \mathbf{C} = \emptyset$, in light of Definition \ref{def:interventional_setting}, this is the same as 
\[\left( (.5-b_s)^{\mid \mathbf{C}\mid}   (.5+b_s)^{\mid \mathbf{X}\setminus \mathbf{C}\mid} (.5+b_w)^{\mid \mathbf{W}\mid} (.5-b_w)^{\mid \mathbf{V}\setminus \mathbf{W}\mid} )  \right),\] which on the assumption of regularity $>0$.

\item Next, we are looking at $P_j(\mathbf{u}, \mathbf{v})$. As $\mathbf{v}$ are actual values from $(M,\mathbf{u})$,  $\mathbbm{1}((M, \mathbf{u}) \vDash \mathbf{v})=1$ and so by \eqref{eq:joint}  $P_j(\mathbf{u}, \mathbf{v}) = P(\mathbf{u})$. Since $(M, \mathbf{u})$ are realized by the assumption, $P(\mathbf{u}) > 0$ and therefore also $P_j(\mathbf{u}, \mathbf{v})>0$. 

\item The value of the indicator $\mathbbm{1}\left(N_{\mathbf{C} = \mathbf{c}, \mathbf{T}}(\mathbf{u}, \varphi)\right)$ is, following \eqref{eq:nec}, determined by there being a $\mathbf{c'}$ such that 
$(M, \mathbf{u}) \vDash [\mathbf{C} \leftarrow \mathbf{c'}, \mathbf{T} \leftarrow \mathbf{t}^\star]\neg \varphi$, which is satisfied by the context-sensitive necessity assumption with the frozen variables as witnesses for the existential claim, so the value of the indicator is 1.

\item For $\mathbbm{1}\left(S_{\mathbf{C} = \mathbf{c}, \mathbf{T}}(\mathbf{u}, \varphi)\right)$,  note that with exogenous noise $\mathbf{u}$ held constant, given the first part of the factivity assumption:
\[\left( (M, \mathbf{u}) 
    \vDash [\mathbf{C} \leftarrow \mathbf{c}, \mathbf{T} \leftarrow \mathbf{t}^\star]\varphi \right) \Leftrightarrow \left( (M, \mathbf{u})\models \varphi\right) \] \noindent  Since the right-hand side holds by the second factivity assumption, so does the left-hand side, which yields $\mathbbm{1}\left(S_{\mathbf{C} = \mathbf{c}, \mathbf{T}}(\mathbf{u}, \varphi)\right) =1$ by \linebreak ($S_{\mathbf{C} = \mathbf{c}, \mathbf{T}}(\mathbf{u}, \varphi)$).

\end{itemize}
\end{proof}




\begin{proof}[Proof of Theorem \ref{th:exp-ac}] Let us fix notation as in the proof of Theorem \ref{th:ac-exp}.  Since the probabilities are non-negative, for the assumption to hold, there has to be at least one $T$ for which \eqref{eq:pect} holds, and so we are now looking again at the decomposition in \eqref{eq:decomposition}. However, now, by the assumption, the whole composite is positive, and therefore so is $\mathbbm{1}\left(N_{\mathbf{C} = \mathbf{c}, \mathbf{T}}(\mathbf{u}, \varphi)\right)$, which by \eqref{eq:nec} means that there are $\mathbf{c'}$ and $\mathbf{T}$ such that: 
\[
    (M, \mathbf{u}) \vDash [\mathbf{C} \leftarrow \mathbf{c'}, \mathbf{T} \leftarrow \mathbf{t}^\star]\neg \varphi.\] These are the witnesses needed to validate the existential claim involved in Definition \eqref{def:ac}.
\end{proof}





\begin{proof}[Proof of Theorem \ref{th:local_max_to_ac}]
We already assumed factivity, and by Theorem \ref{th:if_pr_then_nec} we know the context-sensitive necessity condition is satisfied. So the only reason why $\mathbf{C^\dagger}$ could fail to be an actual cause is minimality. We will assume minimality indeed fails and reason to the conclusion that in that case $(\mathbf{C^\dagger}, \mathbf{T^\dagger})$ couldn't belong to $\Gamma$ in the first place.


So suppose there is a $\mathbf{C^\sharp} \subset \mathbf{C^\dagger}$ which is an actual cause with a witness $\mathbf{T^\sharp}$. $\Phi(\mathbf{C^\dagger}, \mathbf{T^\dagger})$ and $\Phi(\mathbf{C^\sharp}, \mathbf{T^\sharp})$ are, respectively, as follows:
\begin{align}
P^{\mathit{set}}_{b_s,b_w, \{\mathbf{C^\dagger}=\mathbf{c^\dagger}\}, \mathbf{V}}(\mathbf{C^\dagger},\mathbf{T^\dagger}) \cdot P_j(\mathbf{u}, \mathbf{v}) \cdot \mathbbm{1}\left(N_{\mathbf{C^\dagger} = \mathbf{c^\dagger}, \mathbf{T^\dagger}}(\mathbf{u}, \varphi)\right) \cdot \mathbbm{1}\left(S_{\mathbf{C^\dagger} = \mathbf{c^\dagger}, \mathbf{T^\dagger}}(\mathbf{u}, \varphi)\right)  \label{eq:comparisondagger}\\
P^{\mathit{set}}_{b_s,b_w, \{\mathbf{C^\sharp}=\vec {c^\sharp}\}, \mathbf{V}}(\mathbf{C^\sharp},\mathbf{T^\sharp}) \cdot P_j(\mathbf{u}, \mathbf{v}) \cdot \mathbbm{1}\left(N_{\mathbf{C^\sharp} = \mathbf{c^\sharp}, \mathbf{T^\sharp}}(\mathbf{u}, \varphi)\right) \cdot \mathbbm{1}\left(S_{\mathbf{C^\sharp} = \mathbf{c^\sharp}, \mathbf{T^\sharp}}(\mathbf{u}, \varphi)\right) \label{eq:comparisonsharp}
\end{align}

\noindent In comparing them, we can drop  $P_j(\mathbf{u}, \mathbf{v})$ as identical. Moreover, indicators in \eqref{eq:comparisondagger} are 1s on the assumption that $\Phi(\mathbf{C^\dagger}, \mathbf{T^\dagger})>0$ and indicators in \eqref{eq:comparisonsharp} are 1s by Theorem \ref{th:ac-exp} and the assumption that $C^\sharp$ is an actual cause with witness $\mathbf{T^\sharp}$.  So all that is left to compare is $P^{\mathit{set}}_{b_s,b_w, \{\mathbf{S}=\mathbf{s}\}, \mathbf{V}}(\mathbf{C^\dagger},\mathbf{T^\dagger})$ and  $P^{\mathit{set}}_{b_s,b_w, \{\mathbf{S}=\mathbf{s}\}, \mathbf{V}}(\mathbf{C^\sharp},\mathbf{T^\sharp})$.

\begin{align}
P^{\mathit{set}}_{b_s,b_w, \{\mathbf{S}  =\mathbf{s}\}, \mathbf{V}}(\mathbf{C^\dagger},\mathbf{T^\dagger}) & =  P_{b_s, b_w}\left(I_{\mathbf{C^\dagger}} \wedge I^-_{\mathbf{S} \backslash \mathbf{C^\dagger}}\right) \cdot \underbrace{P_{a,w} \left( J_{\mathbf{T^\dagger}}  \wedge J^-_{\mathbf{V} \backslash \mathbf{T^\dagger}} \right)}_{ = .5^{\vert \mathbf{V} \vert} \text{ by $w=0$}}\\
P^{\mathit{set}}_{b_s,b_w, \{\mathbf{S}=\mathbf{s}\}, \mathbf{V}}(\mathbf{C^\sharp},\mathbf{T^\sharp}) & =  P_{b_s, b_w}\left(I_{\mathbf{C^\sharp}} \wedge I^-_{\mathbf{S} \backslash \mathbf{C^\sharp}}  \right) \cdot \overbrace{P_{b_s,b_w}  \left(  J_{\mathbf{T^\sharp}}  \wedge J^-_{\mathbf{V} \backslash \mathbf{T^\sharp}} \right)}
\end{align}
\noindent We are now left comparing $(.5-p_s)^{\vert \mathbf{C^\dagger}\vert } (.5+p_s)^{\vert \mathbf{S}\setminus{\mathbf{C^\dagger}}\vert}$ with $(.5-p_s)^{\vert \mathbf{C^\sharp}\vert } (.5+p_s)^{\vert \mathbf{S}\setminus{\mathbf{C^\sharp}}\vert}$. But by the assumption that $\mathbf{C^\sharp} \subset \mathbf{C^\dagger}$, we know $\vert \mathbf{C^\dagger}\vert > \vert \mathbf{C^\sharp}\vert$ and $\vert \mathbf{S}\setminus{\mathbf{C^\dagger}}\vert < \vert \mathbf{S}\setminus{\mathbf{C^\sharp}}\vert$. Thus $P^{\mathit{set}}_{b_s,b_w, \{\mathbf{S}=\mathbf{s}\}, \mathbf{V}}(\mathbf{C^\sharp},\mathbf{T^\sharp}) > P^{\mathit{set}}_{b_s,b_w, \{\mathbf{S}  =\mathbf{s}\}, \mathbf{V}}(\mathbf{C^\dagger},\mathbf{T^\dagger})$ and so $\Phi(\mathbf{C^\sharp}, \mathbf{T^\sharp}) > \Phi(\mathbf{C^\dagger}, \mathbf{T^\dagger})$, which  contradicts the assumption that $(\mathbf{C^\dagger}, \mathbf{T^\dagger}) \in \Gamma$.
\end{proof}




\begin{proof}[Proof of Corollary \ref{cor:marginalize}]
The assumptions entail there is a $\mathbf{T^\dagger}\in \mathcal{P}(\mathbf{V})$ such that $(\mathbf{C^\dagger}, \mathbf{T^\dagger}) \in \Gamma$ (as defined before) and $\Phi(\mathbf{C^\dagger}, \mathbf{T^\dagger})>0$, because $\{(\mathbf{C^\dagger}, \mathbf{T})\vert \mathbf{T}\subseteq \mathbf{V}\} \supseteq \Gamma$. This gives a witness set that together with $\mathbf{C^\dagger}$ satisfies the conditions of Theorem \ref{th:local_max_to_ac}, so by that theorem it follows that $C^\dagger$ is an actual cause.
\end{proof}


\begin{proof}[Argument for Observation \ref{obs:misalignment}]
Consider a very simple model with two parent binary nodes $A, B$ and a binary child node $Y$, with uniform distribution over the states of $A, B$. Let $C  =  \mathbbm{1}\left( (M, \mathbf{u}) \models A \vee B\right)$. Consider $C^\dagger = \{A = 1, B =1\}$ and $C^{\sharp} = \{A=1\}$. Let $b_w$ be 0, so that all witness activations are equally probable. Let $u_{ij}$ represent a setting where $A=i, B= j$.
\begin{itemize}
\item In $u_{11}$, $C^\dagger$ is an actual cause of $\varphi := Y=1$. Setting the values of $A, B$ to the only alternative results in $\neg\varphi$, no change in any of the subsets to $C^\dagger$ does so, and necessity is satisfied. 
\item There is effectively only one non-zero-measure witness set for $C^\dagger$, which is $\emptyset$. With that witness set, necessity indicators and the sufficiency indicators are always 1 (always, setting $A$ and $b$ to ones results in $Y=1$, and always, setting them to zeros results in $Y=0$. The probability of each scenario is $1/4$ (there are four such settings), and the probability of that causal setting is $(.5-p_s)^2$. So, $P(C^\dagger \rightarrow \varphi) = 4 \times ((.5-p_s)^2 \times 1/4)$.


\item in $u_{10}$, $C^\sharp$ is an actual cause of $\varphi$. There are two witness sets, $\emptyset$ and $\{B\}$, each equally likely. The sufficiency indicator for $C^\sharp$ is 1 everywhere, but the necessity indicator for $C^\sharp$ is 1 in $u_{10}$ and $u_{00}$,  so we only sum over those. So we multiply by 2 because of 2 witness sets, and by 2 again, because there are two settings of $\mathbf{u}$ where both indicators are 1.
This means we multiply by 4 as before. However, this time the probability of the causal setting is $(.5-p_s)(.5+p_s)$, so together we have that $P(C^\sharp \rightarrow \varphi) = 4( (.5-p_s)(.5+p_s) \times 1/4) $
\end{itemize}
This means that the difference between these two probabilities of causation will boil down to differences resulting from the choice of $p_a$. On rare occasions, they will be identical, but for many choices they will not, and only one of $C^\dagger$, $C^\sharp$ will end up in $\Xi$ (see Figure \ref{fig:simple_max}).

\begin{figure}
    \centering
    \includegraphics[width=0.8\linewidth]{figures/set_probabilities.png}
    \caption{Probabilities of causation in a simple three-node model; for each set, there exists a setting $\mathbf{u}$ in which that set is an actual cause of $Y=1$, but whether the probability of causation $P(Set \rightarrow Y=1)$ is maximized depends on the choice of the suspect preemption bias. This ambiguity does disappear once we identify which $\mathbf{u}$ are actual. But this also suggests that inspecting only the maximizers of the probability of causation is not a robust strategy for identifying actual causes across a range of possible values of the exogenous nodes.}
    \label{fig:simple_max}
\end{figure}


\end{proof}




\end{document}

